\documentclass[12pt]{article}
\usepackage{fullpage}
\usepackage{amsmath,amssymb,amsthm}
\usepackage[T1]{fontenc}
\usepackage{hyperref}

\newcommand{\Pone}{\mathbb P^1}
\newcommand{\Sph}{\mathbb S^2}
\newcommand{\C}{\mathbb C}
\newcommand{\Z}{\mathbb Z}
\newcommand{\Teich}{\mathcal T}
\newcommand{\Mod}{\mathcal M}
\newcommand{\D}{\mathcal D}
\newcommand{\W}{\mathcal W}
\newcommand{\V}{\mathcal V}
\newcommand{\wind}{\operatorname{wind}}
\newcommand{\id}{\operatorname{id}}
\newcommand{\Push}{\operatorname{Push}}
\newcommand{\PMod}{\operatorname{PMod}}
\newcommand{\pr}{\operatorname{pr}}
\newtheorem{theorem}{Theorem}[section]
\newtheorem{lemma}[theorem]{Lemma}
\newtheorem{proposition}[theorem]{Proposition}
\newtheorem{construction}[theorem]{Construction}
\theoremstyle{remark}
\newtheorem{remark}[theorem]{Remark}
\newtheorem{convention}[theorem]{Convention}
\numberwithin{equation}{section}

\title{Problem 7 -- ProofCouncil}
\author{}
\date{}

\begin{document}
\maketitle

\begin{abstract}
Let $f$ be a quadratic rational map one of whose critical points lies in a
cycle $A$ of length four and whose other critical value $v$ lies outside
$A$, and let $B=A\cup\{v\}$. Epstein's deformation space $\D_f\subset\Teich(\Sph,B)$ is
known to be disconnected (Hironaka--Koch). We write down an explicit simple
closed curve $\gamma$ on $(\Sph,B)$, for a normalized model map the image
of a round circle under two Dehn twists about a second round circle,
transported to a general $f$ along a path in parameter space, and prove
that the Dehn twist $T_\gamma$ itself preserves $\D_f$ and that the
translates $T_\gamma^n(D_0)$, $n\in\Z$, of the base component $D_0$ are
pairwise distinct components, using the covering-space framework of
Hironaka--Koch and a winding-number obstruction.
\end{abstract}

\section{Introduction}

\subsection{Deformation spaces}\label{sec:defspaces}

Throughout, $\Pone$ is the Riemann sphere, which we also regard as the
oriented topological sphere $\Sph$. For a finite set $B\subset\Pone$ with
$|B|\ge3$, the \emph{Teichm\"uller space} $\Teich(B)=\Teich(\Sph,B)$ is the
set of orientation-preserving homeomorphisms $\phi\colon\Pone\to\Pone$
(\emph{markings}), where $\phi\sim\phi'$ when $\phi'=\mu\circ\phi\circ h$
for a M\"obius transformation $\mu$ and a homeomorphism $h$ isotopic to the
identity relative to $B$. Pulling back the standard complex structure by
$\phi$ identifies $\Teich(B)$ with the complex structures on $(\Sph,B)$ up
to isotopy relative to $B$; the class $\tau_0=[\id]$ is the \emph{base
point}.

Let $f\colon\Pone\to\Pone$ be an orientation-preserving branched cover of
degree $d\ge2$ with set of critical values $V(f)$, and let $A,B\subset\Pone$
be finite sets with $|A|,|B|\ge3$ and $f(A)\cup V(f)\subset B$.
Pulling back the complex structure of a marking $\phi$ under $f$ gives a
complex structure on $(\Sph,A)$; concretely, there are a homeomorphism
$\psi$, unique up to post-composition with a M\"obius transformation, and a
rational map $F$ with
\begin{equation}\label{eq:pullbackdiagram}
 F\circ\psi=\phi\circ f ,
\end{equation}
and $f^*[\phi]:=[\psi]\in\Teich(A)$. This is the Thurston pullback map
$f^*\colon\Teich(B)\to\Teich(A)$; it is holomorphic
\cite[\S1.3, Prop.\ 1.4]{BCT}. We call $\phi$ the \emph{target marking} and
$\psi$ the \emph{source marking} of the diagram
\eqref{eq:pullbackdiagram}. If $A\subset B$, forgetting the points of
$B\setminus A$ gives the holomorphic \emph{forgetful map}
$i_*\colon\Teich(B)\to\Teich(A)$, $[\phi]\mapsto[\phi]$. When
$f(A)\cup V(f)\cup A\subset B$, both maps are defined, and the
\emph{deformation space} of $f$ is their equalizer
\[
 \D_f=\{\tau\in\Teich(B):f^*(\tau)=i_*(\tau)\}.
\]
This notion is due to Epstein \cite{Ep}; see \cite[\S1.4]{BCT} for an
exposition. When $f$ is rational, as in the problem considered here, we
may take $\psi=\id$, $F=f$ in \eqref{eq:pullbackdiagram} for $\phi=\id$,
so $\tau_0\in\D_f$. (By Epstein's theorem \cite[Thm.\ 1.5]{BCT}, under the
standing hypotheses above, if $f$ is not a topological $(2,2,2,2)$-map in
the sense of \cite[\S1.4]{BCT}, then $\D_f$, whenever nonempty, is a smooth
complex submanifold of $\Teich(B)$ of dimension $|B\setminus A|$. For the
rational family considered here $\tau_0\in\D_f$, and the periodic critical
point excludes the exceptional type.)

The \emph{pure mapping class group} $\PMod(B)$, the group of isotopy classes
relative to $B$ of orientation-preserving homeomorphisms of $\Pone$ fixing
$B$ pointwise, acts on $\Teich(B)$ on the right by
\begin{equation}\label{eq:rightaction}
 [\phi]\cdot h=[\phi\circ h];
\end{equation}
for a subset $D\subset\Teich(B)$ we write $h(D)=\{\tau\cdot h:\tau\in D\}$.
A mapping class that preserves $\D_f$ permutes its connected components.
For a simple closed curve $\gamma$ in $\Pone\setminus B$ we write
$T_\gamma\in\PMod(B)$ for the Dehn twist about $\gamma$. The question
addressed here asks for an explicit $\gamma$ such that some power of
$T_\gamma$ preserves $\D_f$ and has an infinite orbit on the components.

\subsection{The family and the result}\label{sec:family}

Let $f$ be a quadratic rational map with critical points $c_1,c_2$ such that
$c_1$ lies in a cycle of length $4$ and the critical value $v=f(c_2)$ is not
in this cycle, and put
\[
 A=\{c_1,f(c_1),f^{\circ2}(c_1),f^{\circ3}(c_1)\},\qquad B=A\cup\{v\}.
\]
Then $f(A)=A$ and $V(f)=\{f(c_1),v\}\subset B$, so $\D_f\subset\Teich(B)$ is
defined; $\Teich(B)$ has dimension $2$. Conjugating $f$ by the M\"obius
transformation that sends $c_1,f(c_1),f^{\circ2}(c_1)$ to $0,\infty,1$
normalizes the cycle to $0\mapsto\infty\mapsto1\mapsto s\mapsto0$, and
Section~\ref{sec:curve} shows that this forces $f$ to become
\begin{equation}\label{eq:family}
 F_s(z)=1+\frac{-s^2-s+1}{(s-1)z}+\frac{s(2s-1)}{(s-1)z^2},
 \qquad s\in U,
\end{equation}
for an explicit Zariski-open set $U\subset\C$ given in \eqref{eq:U}. Our
model is $F_a$ with $a=3/5$, whose marked set is
$B_a=\{0,\infty,1,a,v_a\}$, where $v_a=121/120$ is its second critical
value (computed in \eqref{eq:vs}).

The answer (Theorem~\ref{thm:main}) is as follows. Let $C$ be the round
circle $|z|=241/240$, which separates $\{\infty,v_a\}$ from $\{0,1,a\}$,
and let $\eta$ be the round circle $|z-4/5|=1/4$, which encloses exactly
$a,1,v_a$. Put $\gamma_a=T_\eta^{2}(C)$, the image of $C$ under two Dehn
twists about $\eta$ (with the sign of Convention~\ref{conv:sign}; either
sign works, Remark~\ref{rem:othersign}). For a general $f$ of the family,
let $\gamma$ be obtained from $\gamma_a$ by following any path in $U$ from
$a$ to the parameter of $f$ and dragging the marked points along. Then
$T_\gamma$ preserves $\D_f$, and the translates $T_\gamma^{n}(D_0)$,
$n\in\Z$, of the component $D_0\ni\tau_0$ are pairwise distinct
components; so the power of the twist can be taken to be $1$. (A path is
needed: $\gamma_a$ and $C$ separate $B_a$ in the same way, so separation
data do not determine the isotopy class of $\gamma$.)

\subsection{Relation to earlier work}\label{sec:prior}

Hironaka and Koch \cite{HK} proved that for this family $\D_f$ has
infinitely many connected components (\cite[Thm.\ 1.4]{HK}), introducing
the tools we use: the surface of marked quadratic maps with its
degree-four covering of moduli space and its diagonal
(\cite[\S\S4.1--4.2]{HKa}, \cite[\S3]{HK}), the description of $\D_f$
through mapping classes that lift under $f$ and through equalizers
(\cite[\S\S2.8--2.11]{HK}), and a symmetric quotient of the parameter
surface used as an obstruction (\cite[\S2.12]{HK}); they also exhibit a
mapping class preserving $\D_f$ with an infinite orbit on components
(\cite[Prop.\ 3.2]{HK}). Here \cite{HKa} is the first arXiv version of
their paper and \cite{HK} the revised author version, whose numbering we
follow. Firsova, Kahn and Selinger \cite{FKS} treat the same family
(\cite[\S6]{FKS}; numbering of arXiv v4), study twists along multicurves
adapted to the pullback (``equalizing multicurves''), and prove that these
deformation spaces are not contractible (\cite[Lemma 8.4]{FKS}). The
present note adds the identification of a single, explicitly drawn Dehn
twist with an infinite component orbit, which we have not found in
\cite{HK} or \cite{FKS}, with an elementary verification; our obstruction
map $\Phi$ (Section~\ref{sec:detector}) is a form of the symmetric
quotient of \cite{HK}.

\subsection{Strategy and outline}\label{sec:strategy}

The proof has five steps. (1) Section~\ref{sec:covering} factors the
universal cover $\Teich(B_a)\to M$ of the moduli space $M$ of five marked
points through a finite cover $\W\to M$, the surface of quadratic maps with
the prescribed portrait of marked orbits; the pullback and forgetful maps
become the coordinate projections of $\W$, and $\D_{F_a}$ lies over the
diagonal $\V\subset\W$. (2) Section~\ref{sec:criterion}: a loop $g$ in
$\W$ acts on $\Teich(B_a)$ as a mapping class, which preserves $\D_{F_a}$
if the two coordinate loops of $g$ are homotopic, and preserves the base
component $D_0$ if and only if $g$ is homotopic in $\W$ to a loop in $\V$.
(3) Section~\ref{sec:loops} exhibits a loop $g=\ell^{-1}m\ell$ in $\W$
acting as $T_{\gamma_a}^{\pm1}$: $m$ is a meridian acting as the twist
about $C$, and the point push along $\ell$ is identified exactly by a
Joukowski substitution. (4) Section~\ref{sec:detector} shows that no power
$g^n$, $n\ne0$, is homotopic to a loop in $\V$, using a map
$\Phi=(t,L)\colon\W\to(\C\setminus\{0,1\})\times\C^*$ whose first
coordinate becomes, on an enlarged diagonal, the two-sheeted covering
$w\mapsto w^2$, while its second coordinate imposes a linear relation
between two winding numbers that the loops $g^n$ violate. (5)
Section~\ref{sec:transport} transports the conclusion to every map of the
family by lifting an isotopy of marked points through the branched covers.

\section{Conventions on loops, point pushing and twists}\label{sec:conventions}

For $B$ as above, the \emph{moduli space} $\Mod(B)$ is the set of injective
maps $B\to\Pone$ modulo post-composition with M\"obius transformations.
Choosing three points of $B$ and normalizing their images to $0,\infty,1$
identifies $\Mod(B)$ with an open subset of $\C^{|B|-3}$. The projection
$\pi_B\colon\Teich(B)\to\Mod(B)$, $[\phi]\mapsto[\phi|_B]$, is a holomorphic
universal covering whose deck group is $\PMod(B)$ acting by
\eqref{eq:rightaction}; see \cite[\S2.1]{KPS}.

Two moduli spaces occur. For four marked points labelled $0,\infty,1,x$,
and for five points labelled $0,\infty,1,y,z$, we use respectively
\begin{equation}\label{eq:YM}
 Y=\C\setminus\{0,1\},
 \qquad
 M=\{(y,z)\in Y^2:y\ne z\}.
\end{equation}
Thus for $A\supset\{0,\infty,1\}$ with $|A|=4$, $\pi_A\colon\Teich(A)\to Y$
records the fourth point under the marking normalized to fix $0,\infty,1$,
and similarly $\pi_B\colon\Teich(B)\to M$ when $|B|=5$.

Products of based loops are read from right to left: $\alpha\beta$
traverses $\beta$ first and then $\alpha$, matching composition of
homeomorphisms.\label{conv:products} For a closed loop $\omega$ in
$\C\setminus\{c\}$ we write
$\wind_c(\omega)\in\Z$ for its winding number about $c$; for a non-closed
path, $\wind_c(\omega)$ denotes the total change of argument of $\omega-c$
divided by $2\pi$, which is additive under concatenation and changes sign
under reversal.

Let $\delta$ be a loop in $\Mod(B)$ based at $\pi_B(\tau_0)$, that is, a
motion of the marked points. By isotopy extension there is an ambient
isotopy $h_t$ of $\Pone$ with $h_0=\id$ such that $h_t|_B$ follows
$\delta(t)$. The mapping class $P(\delta)=[h_1]\in\PMod(B)$ depends only on
the homotopy class of $\delta$, and $[h_t]$ is the lift of $\delta$ to
$\Teich(B)$ starting at $\tau_0$, with endpoint $\tau_0\cdot P(\delta)$.
Consequently the deck transformation associated with $\delta$, the one
sending $\tau_0$ to the endpoint of this lift, is
\begin{equation}\label{eq:deckaction}
 [\phi]\longmapsto[\phi\circ P(\delta)].
\end{equation}
Point pushing respects products, $P(\alpha\beta)=P(\alpha)P(\beta)$, since
the isotopy following $\alpha\beta$ may be taken to be $h^\beta_t$ followed
by $h^\alpha_t\circ h^\beta_1$. Because \eqref{eq:deckaction} is a right
action, deck transformations of products compose in the opposite order;
we only use deck transformations of powers of a single loop. We write
$g(\tau)$ for the deck transformation \eqref{eq:deckaction} of a loop $g$
applied to $\tau$.

When only one marked point $v$ moves, $\delta$ is a loop $\zeta$ in
$\Pone\setminus(B\setminus\{v\})$ based at $v$, and $P(\zeta)$ is the
\emph{point-pushing map} $\Push(\zeta)$ of the Birman exact sequence
\cite[\S4.2, Thm.\ 4.6]{FM}. We use two standard facts.
(P1) If $\zeta$ is a simple loop, then
$\Push(\zeta)=T_{\zeta_1}T_{\zeta_2}^{-1}$, where $\zeta_1,\zeta_2$ are the
two boundary curves of a regular neighbourhood of $\zeta$, the order being
fixed by the orientation of $\zeta$ and the sign convention for twists
\cite[Fact 4.7]{FM}.
(P2) For every mapping class $h$ and simple closed curve $c$,
$hT_ch^{-1}=T_{h(c)}$ \cite[Fact 3.7]{FM}.

Two consequences of (P1) are used. First, if $\Delta\subset\Pone$ is a
closed disk with $\partial\Delta\cap B=\emptyset$ and
$\Delta\cap B=\{v,w\}$, and $\zeta$ is a simple loop of $v$ in
$\Delta\setminus\{w\}$ going once around $w$, then one boundary curve of
its regular neighbourhood bounds a disk containing only $w$ (trivial
twist) and the other is isotopic to $\partial\Delta$, so
\begin{equation}\label{eq:twopunctures}
 \Push(\zeta)=T_{\partial\Delta}^{\pm1}.
\end{equation}
Second, if $\kappa,\eta$ are disjoint simple closed curves bounding an
annulus meeting $B$ only in $v$, with $\kappa$ inside $\eta$ in a given
disk picture, and $\zeta$ is a simple loop of $v$ once counterclockwise
around the core of the annulus, then
$\Push(\zeta)=T_\eta^{\epsilon}T_\kappa^{-\epsilon}$ with $\epsilon=\pm1$
depending only on the sign convention for Dehn twists.

\begin{convention}\label{conv:sign}
We fix the sign convention for Dehn twists so that $\epsilon=+1$, that is,
$\Push(\zeta)=T_\eta T_\kappa^{-1}$ for such a counterclockwise circuit.
With the opposite convention the boxed formula \eqref{eq:answer} of
Section~\ref{sec:curve} reads $T_\eta^{-2}(C)$; both curves have the
stated property (Remark~\ref{rem:othersign}).
\end{convention}

\section{The model map and the curve}\label{sec:curve}

Let $f$ be as in Section~\ref{sec:family} and conjugate it so that its
critical cycle is $0\mapsto\infty\mapsto1\mapsto s\mapsto0$ with $0$ the
critical point. A quadratic rational map with a critical point at $0$
mapping to $\infty$ has a double pole at $0$ and no other pole, so it is
$1+b_s/z+c_s/z^2$ after using $F(\infty)=1$. Solving the linear equations
$F(1)=s$ and $F(s)=0$ for $b_s,c_s$ gives $b_s=(-s^2-s+1)/(s-1)$ and
$c_s=s(2s-1)/(s-1)$, which is \eqref{eq:family}. The other critical point of $F_s$ is $-2c_s/b_s$, and
its critical value is $1-b_s^2/(4c_s)$, which simplifies to
\begin{equation}\label{eq:vs}
 v(s)=-\frac{(s^2-3s+1)^2}{4s(s-1)(2s-1)} .
\end{equation}
The map has degree two and the four cycle points are distinct exactly when
$s\ne0,1,1/2$ (for $s=1/2$ one has $c_s=0$). The identities
\[
 v(s)-1=-\frac{(s^2+s-1)^2}{4s(s-1)(2s-1)},\qquad
 v(s)-s=-\frac{(3s^2-3s+1)^2}{4s(s-1)(2s-1)}
\]
show that $v(s)\notin\{0,\infty,1,s\}$, that is, $v(s)$ lies outside the
cycle, exactly when the three quadratic factors below are nonzero. So the
maps in the problem are, up to M\"obius conjugacy, precisely the $F_s$ with
\begin{equation}\label{eq:U}
 s\in U=\left\{s\in\C:
 s(s-1)(2s-1)(s^2-3s+1)(s^2+s-1)(3s^2-3s+1)\ne0\right\}.
\end{equation}
The set $U$, a cofinite subset of $\C$, is connected. We write
$A_s=\{0,\infty,1,s\}$ and $B_s=A_s\cup\{v(s)\}$ for the marked sets of
$F_s$. Our model is the parameter $a=3/5$:
\begin{equation}\label{eq:model}
 F_a(z)=1-\frac1{10z}-\frac3{10z^2},\qquad
 0\longmapsto\infty\longmapsto1\longmapsto a\longmapsto0,
\end{equation}
with other critical point $-6$ and other critical value
$v_a=v(3/5)=121/120$. On the sphere marked by
$B_a=\{0,\infty,1,a,v_a\}$ draw the round circles
\begin{equation}\label{eq:circles}
 C=\left\{|z|=\tfrac{241}{240}\right\},\qquad
 \eta=\left\{\left|z-\tfrac45\right|=\tfrac14\right\}.
\end{equation}
The exterior of $C$ contains exactly the marked points $\infty,v_a$
(note $v_a=242/240$), and the interior of $\eta$ contains exactly $a,1,v_a$
(the distances from $4/5$ to $a$, $1$, $v_a$ are $1/5$, $1/5$, $25/120$,
all less than $1/4$, while $|0-4/5|>1/4$). The two circles intersect. Our
model curve is
\begin{equation}\label{eq:answer}
 \boxed{\ \gamma_a=T_\eta^{\,2}(C),\ }
\end{equation}
with the sign of Convention~\ref{conv:sign}.

\begin{construction}[Transport to a general map]\label{con:transport}
Let $f$ be a map of the family with normalized form $F_s$, $s\in U$, the
normalizing M\"obius transformation being $\mu$ (so $\mu f\mu^{-1}=F_s$).
Choose a path $s_t$, $0\le t\le1$, in $U$ from $a$ to $s$. The five points
$0,\infty,1,s_t,v(s_t)$ stay distinct, so by isotopy extension there is an
ambient isotopy $H_t$ of $\Pone$ with $H_0=\id$ fixing $0,\infty,1$ and
with $H_t(a)=s_t$, $H_t(v_a)=v(s_t)$. Put
$\gamma=\mu^{-1}(H_1(\gamma_a))\subset\Pone\setminus B$; equivalently,
$H_1(\gamma_a)$ is the image of $H_1(C)$ under two twists about $H_1(\eta)$.
\end{construction}

The isotopy class of $\gamma$ relative to $B$ depends on the path $s_t$
only through its homotopy class in $U$; Theorem~\ref{thm:main} holds for
every choice.

\begin{theorem}\label{thm:main}
Let $f$, $A$, $B$ be as in Section~\ref{sec:family}, and let
$\gamma$ be the curve of Construction~\ref{con:transport}. Then the Dehn
twist $T_\gamma$ preserves $\D_f$. If $D_0$ denotes the connected component
of $\D_f$ containing the base point $\tau_0$, then the components
$T_\gamma^{n}(D_0)$, $n\in\Z$, are pairwise distinct.
\end{theorem}

Sections \ref{sec:covering}--\ref{sec:detector} prove the theorem for the
model $f=F_a$ and $\gamma=\gamma_a$; Section~\ref{sec:transport} deduces
the general case.

\section{A finite cover of the moduli space}\label{sec:covering}

Recall $Y$ and $M$ from \eqref{eq:YM}; $M$ is the moduli space of
$B_a$, with $0,\infty,1$ normalized and $(y,z)$ the positions of the points
corresponding to $a$ and $v_a$. A point of $\Teich(B_a)$ determines,
through \eqref{eq:pullbackdiagram}, a quadratic rational map with marked
source and target points; the normalized positions of the fourth source and
target points give a map from $\Teich(B_a)$ to a surface $\W$ of marked
maps, a finite cover of $M$. This is the parameter surface of Hironaka and
Koch \cite[\S\S4.1--4.2]{HKa}; in general such spaces are Hurwitz spaces,
finite coverings of the moduli space of the target \cite[Thm.\ 2.6]{Ko},
\cite[\S3.1]{KPS}.

Consider a quadratic rational map with marked source points $0,\infty,1,x$
and marked target points $0,\infty,1,y,z$ realizing the portrait
\begin{equation}\label{eq:portrait}
 0\stackrel{2}{\longmapsto}\infty,\qquad
 \infty\longmapsto1,\qquad 1\longmapsto y,\qquad
 x\longmapsto0,
\end{equation}
where the $2$ indicates a critical point, and having $z$ as its other
critical value. A degree-two map with a double pole at $0$ and value $1$ at
$\infty$ is $(t-x)(t-r)/t^2$ with zeros $x$ and $r$, and the condition
$1\mapsto y$ determines $r$. Hence the map is unique:
\begin{equation}\label{eq:Fxy}
 F_{x,y}(t)=\frac{(t-x)(t-r)}{t^2},
 \qquad r=\frac{x+y-1}{x-1}.
\end{equation}
Its critical points are $0$ and $2xr/(x+r)$, and the critical value at the
latter is $-(x-r)^2/(4xr)$, which in terms of $x,y$ reads
\begin{equation}\label{eq:Z}
 Z(x,y)=\frac{((x-1)^2-y)^2}{4x(x+y-1)(1-x)} .
\end{equation}
The marked points must be distinct and the map must have degree two
($r\ne0$), so the parameter surface and its projection to $M$ are
\begin{align}
 \W&=\{(x,y)\in Y^2:x+y\ne1,\ Z(x,y)\notin\{0,1,y\}\},
 \label{eq:W}\\
 p\colon\W&\longrightarrow M,\qquad p(x,y)=(y,Z(x,y)).
 \label{eq:p}
\end{align}
The maps $F_s$ of the family \eqref{eq:family} are the marked maps whose
fourth source and target points agree: indeed $F_{s,s}=F_s$, the two
expressions being quadratic maps with the same poles, zeros and value at
$\infty$, and $Z(s,s)=v(s)$. The conditions $Z(s,s)\ne0,1,s$ are precisely
the three quadratic factors in \eqref{eq:U}, so the \emph{diagonal} of
$\W$ is
\begin{equation}\label{eq:V}
 \V=\W\cap\{x=y\}=\{(s,s):s\in U\},
\end{equation}
a connected Riemann surface, since $U$ is connected.

\begin{lemma}\label{lem:cover}
\begin{enumerate}
\item The map $p$ is an unramified covering of degree four, and $\W$ is
connected.
\item For the model \eqref{eq:model}, with $\pi_A\colon\Teich(A_a)\to Y$
and $\pi_B\colon\Teich(B_a)\to M$ the projections of
Section~\ref{sec:conventions}, the map
\[
 q=(\pi_A\circ F_a^*,\ \pi_A\circ i_*)\colon\Teich(B_a)\longrightarrow\W
\]
is well defined and satisfies $p\circ q=\pi_B$ and $q(\tau_0)=(a,a)$, and
it is the universal covering of $\W$. Thus $\pi_B$ factors through $\W$,
and, with $\pr_x,\pr_y$ the two coordinate projections of $\W$, by
construction
\begin{equation}\label{eq:liftmaps}
 \pi_A\circ F_a^*=\pr_x\circ q,\qquad
 \pi_A\circ i_*=\pr_y\circ q .
\end{equation}
\end{enumerate}
\end{lemma}

\begin{proof}
(1) Fix $(y,z)\in M$. The equation $Z(x,y)=z$ is
\[
 ((x-1)^2-y)^2+4zx(x-1)(x+y-1)=0,
\]
a monic polynomial equation of degree four in $x$. Its values at
$x=0,1,1-y$ are $(1-y)^2$, $y^2$ and $y^2(y-1)^2$, all nonzero, so every
root satisfies $x\in Y$ and $x+y\ne1$; since $z\notin\{0,1,y\}$, every root
gives a point of $\W$. Thus $p^{-1}(y,z)$ consists of the roots of this
equation. The roots are distinct because of the factorizations
\begin{align*}
 Z-1&=-\frac{(x^2+y-1)^2}{4x(x-1)(x+y-1)},\\
 Z-y&=-\frac{(x^2+2xy-2x-y+1)^2}{4x(x-1)(x+y-1)},\\
 \frac{\partial Z}{\partial x}
 &=-\frac{(x^2+y-1)\,((x-1)^2-y)\,(x^2+2xy-2x-y+1)}
 {4x^2(x-1)^2(x+y-1)^2},
\end{align*}
which show that $\partial Z/\partial x$ vanishes only where $Z=1$, $Z=0$
or $Z=y$, and these values of $z$ are excluded in $M$. So $p$ has exactly
four preimages over each point, $\partial Z/\partial x\ne0$ on $\W$, and
the implicit function theorem shows that $p$ is a local homeomorphism. A
local homeomorphism between Hausdorff spaces all of whose fibres have the
same finite cardinality is a covering (intersect the images of disjoint
neighbourhoods of the four preimages of a point, and count), so $p$ is a
covering of degree four. Finally $\W$ is the complement in $\C^2$ of
the finitely many algebraic curves listed in \eqref{eq:W}, and such a
complement is path connected: join two points by
two complex lines through a generic auxiliary point, each meeting the
curves in finitely many points, which a path in the line can avoid.

(2) Let $\tau=[\phi]\in\Teich(B_a)$, with $\phi$ normalized to fix
$0,\infty,1$, and let $F\circ\psi=\phi\circ F_a$ be a pullback diagram
\eqref{eq:pullbackdiagram} with $\psi$ also normalized to fix $0,\infty,1$
(possible since $\psi$ is unique up to a M\"obius transformation). Put
$x=\psi(a)$, $y=\phi(a)$, $z=\phi(v_a)$. Applying the diagram to the
points of $A_a$ and to the critical points $0,-6$ of $F_a$ shows that $F$
has the portrait \eqref{eq:portrait} with marked source points
$\{0,\infty,1,x\}$, marked target points $\{0,\infty,1,y,z\}$, and other
critical value $F(\psi(-6))=z$. Hence $F=F_{x,y}$ and $Z(x,y)=z$ by the
uniqueness statement above; $x,y\in Y$ and $z\notin\{0,1,y\}$ because
$\psi,\phi$ are injective, and $x+y\ne1$ because $F$ has degree two.
Therefore $(x,y)=(\pi_A F_a^*\tau,\pi_A i_*\tau)\in\W$ and
$p(x,y)=(y,z)=\pi_B(\tau)$. For $\tau_0$ we may take $\psi=\phi=\id$, so
$q(\tau_0)=(a,a)$; and $q$ is continuous because $F_a^*$, $i_*$, $\pi_A$
are.

It remains to see that $q$ is a covering. Let $N\subset M$ be an open ball
evenly covered by both $p$ and $\pi_B$, with $p^{-1}(N)=\bigsqcup_jN_j$ and
$\pi_B^{-1}(N)=\bigsqcup_k\widetilde N_k$. The connected set
$q(\widetilde N_k)$ lies in a single $N_j$, where
$q|_{\widetilde N_k}=(p|_{N_j})^{-1}\circ\pi_B|_{\widetilde N_k}$ is a
homeomorphism onto $N_j$. Hence each $N_j$ is evenly covered by $q$
(possibly with empty preimage), so $q(\Teich(B_a))$ is open and closed in
the connected space $\W$, and $q$ is a surjective covering; as
$\Teich(B_a)$ is simply connected, it is the universal covering.
\end{proof}

\section{When a loop preserves a component}\label{sec:criterion}

We now set up the covering-space criterion of Hironaka and Koch
\cite[\S\S2.8--2.11]{HK} in the coordinates of Lemma~\ref{lem:cover}.
Let $b=(a,a)=q(\tau_0)$. A loop $g$ in $\W$ based at $b$ determines, via
the loop $p\circ g$ in $M$ and \eqref{eq:deckaction}, a deck
transformation $\delta$ of $\pi_B$, which sends $\tau_0$ to the endpoint
$\tau_1$ of the lift of $g$ to $\Teich(B_a)$ starting at $\tau_0$. It is
also a deck transformation of $q$: $q\circ\delta$ and $q$ are both lifts
of $\pi_B$ through $p$, and they agree at $\tau_0$ because $q(\tau_1)=b$.
We denote $\delta$ by the same letter $g$ and write $g(\tau)$, as in
Section~\ref{sec:conventions}. The two \emph{coordinate loops} of
$g$ are the loops $\pr_x\circ g$ and $\pr_y\circ g$ in $Y$ based at $a$.

\begin{lemma}\label{lem:criterion}
Let $E=\operatorname{im}\bigl(\pi_1(\V,b)\to\pi_1(\W,b)\bigr)$ be the set
of classes of loops in $\W$ that are homotopic in $\W$ to a loop in the
diagonal $\V$. Let $g\in\pi_1(\W,b)$, acting on $\Teich(B_a)$ as above,
and let $D_0$ be the connected component of $\D_{F_a}$ containing
$\tau_0$.
\begin{enumerate}
\item If the two coordinate loops of $g$ have equal classes in
$\pi_1(Y,a)$, then $g(\D_{F_a})=\D_{F_a}$.
\item $g(D_0)=D_0$ if and only if $g\in E$.
\end{enumerate}
Consequently, if the coordinate loops of $g$ are homotopic and
$g^n\notin E$ for every $n\ne0$, then $g$ preserves $\D_{F_a}$ and the
components $g^n(D_0)$, $n\in\Z$, are pairwise distinct.
\end{lemma}

\begin{proof}
(1) By \eqref{eq:liftmaps}, $F_a^*$ and $i_*$ are lifts through the
universal covering $\pi_A\colon\Teich(A_a)\to Y$ of the maps
$\pr_x\circ q$ and $\pr_y\circ q$, and they take the same value at
$\tau_0$ because $\tau_0\in\D_{F_a}$. Since $g$ is a deck transformation of
$q$, the map $F_a^*\circ g$ is another lift of $\pr_x\circ q$, so
$F_a^*\circ g=\delta_x\circ F_a^*$ for a deck transformation $\delta_x$ of
$\pi_A$. To identify $\delta_x$, lift $g$ to a path $\tilde g$ from $\tau_0$
to $g(\tau_0)$; then $F_a^*\circ\tilde g$ is the lift, starting at
$F_a^*(\tau_0)$, of the loop $\pr_x\circ g$, so $\delta_x$ is the deck
transformation of $\pi_A$ determined by the class of $\pr_x\circ g$
relative to the base point $F_a^*(\tau_0)$. Likewise
$i_*\circ g=\delta_y\circ i_*$ with $\delta_y$ determined by the class of
$\pr_y\circ g$ relative to the same base point $i_*(\tau_0)=F_a^*(\tau_0)$.
Equal classes give $\delta_x=\delta_y=:\delta$, and then for
$\tau\in\D_{F_a}$,
\[
 F_a^*(g\tau)=\delta F_a^*(\tau)=\delta\, i_*(\tau)=i_*(g\tau),
\]
so $g(\D_{F_a})\subset\D_{F_a}$. The coordinate loops of $g^{-1}$ are the
inverses of those of $g$, hence also homotopic, so
$g^{-1}(\D_{F_a})\subset\D_{F_a}$ as well.

(2) First, $\D_{F_a}\subset q^{-1}(\V)$: if $F_a^*\tau=i_*\tau$, the two
coordinates of $q(\tau)$ agree by \eqref{eq:liftmaps}. Next, on a connected
component $N$ of $q^{-1}(\V)$, the maps $F_a^*|_N$ and $i_*|_N$ are lifts
through $\pi_A$ of one and the same map $\tau\mapsto\pr_xq(\tau)=\pr_yq(\tau)$,
so by uniqueness of lifts they agree on all of $N$ if they agree at one
point. Hence $\D_{F_a}$ is a union of connected components of
$q^{-1}(\V)$; these are path connected, $q\colon q^{-1}(\V)\to\V$ being a
covering of the connected manifold $\V$, and $D_0$ is the component
through $\tau_0$.

The deck transformation $g$ permutes the components of $q^{-1}(\V)$, so
$g(D_0)=D_0$ if and only if there is a path in $D_0$ from $\tau_0$ to
$g(\tau_0)$. Given such a path, its image under $q$ is a loop in $\V$
based at $b$, homotopic in $\W$ to $g$ because both lift to paths in the
simply connected space $\Teich(B_a)$ with the same endpoints; so $g\in E$.
Conversely, if $g$ is homotopic in $\W$ to a loop $g'$ in $\V$, the lift
of $g'$ starting at $\tau_0$ stays in $q^{-1}(\V)$, hence in $D_0$, and
ends at $g(\tau_0)$ by homotopy lifting; so $g(\tau_0)\in D_0$.

Finally, if $g^j(D_0)=g^k(D_0)$ with $j\ne k$, then $g^{j-k}(D_0)=D_0$, so
$g^{j-k}\in E$ by (2), contrary to hypothesis.
\end{proof}

\section{Loops in \texorpdfstring{$\W$}{W} representing the twist about
\texorpdfstring{$\gamma_a$}{gamma\_a}}\label{sec:loops}

The goal of this section is a loop $g$ in $\W$ with
$P(p\circ g)=T_{\gamma_a}^{\pm1}$ and trivial coordinate loops
(Proposition~\ref{prop:g}). We work in the fibre
$\W_a=\{x\in\C:(x,a)\in\W\}$ over $y=a=3/5$, along which $p$ fixes the
four points of $A_a$ and moves only $v_a$, so that mapping classes are
point pushes of $v_a$. By \eqref{eq:W} and the
factorizations in the proof of Lemma~\ref{lem:cover}, the points of $\C$
excluded from $\W_a$ are
\begin{equation}\label{eq:forbidden}
 0,\ 1,\ \tfrac25,\qquad
 1\pm\sqrt{\tfrac35}\ \ (Z=0),\qquad
 \pm\sqrt{\tfrac25}\ \ (Z=1),\qquad
 \tfrac25\pm\tfrac{i\sqrt6}{5}\ \ (Z=a);
\end{equation}
here $2/5=1-a$ is the pole of $Z(\cdot,a)$. Their moduli are
$0$, $1$, $0.4$, $1.77\ldots$, $0.22\ldots$, and $\sqrt{2/5}=0.632\ldots$
for the last four; none is $a=0.6$, so
$\ell(\theta)=(ae^{2\pi i\theta},a)$, $0\le\theta\le1$, is a loop in $\W$
based at $b$. Its image $p\circ\ell$ in $M$ moves $v_a$ along the loop
$\zeta(\theta)=Z(ae^{2\pi i\theta},a)$ in $\Pone\setminus A_a$ based at
$v_a=Z(a,a)$.

Let $m$ be a positive meridian about the pole $x=2/5$ in $\W_a$: travel
from $a$ along the real interval toward $2/5$, go once counterclockwise
around a small circle centred at $2/5$, and return along the interval.

\begin{lemma}\label{lem:m}
The mapping class $P(p\circ m)$ is $T_C^{\pm1}$, a single Dehn twist about
the round circle $C$ of \eqref{eq:circles}.
\end{lemma}

\begin{proof}
On $(2/5,3/5]$ the factorizations give $\partial Z(x,a)/\partial x<0$,
because in the numerator of $\partial Z/\partial x$ the factors $x^2-2/5$
and $(x-1)^2-3/5$ are negative and $x^2-\frac45x+\frac25$ is positive.
Moreover $Z(a,a)=v_a$, so $Z(x,a)\ge v_a>1$ on $(2/5,a]$ by
monotonicity; the pole of $Z(\cdot,a)$
at $2/5$ is simple, and $Z(x,a)\to+\infty$ as
$x\downarrow2/5$. Thus $p\circ m$ moves $v_a$ out along the positive real
ray, once around $\infty$ (a small circle about a simple pole maps to a
loop winding once about $\infty$), and back. For a small enough circle this motion takes place in the closed disk
$\Delta=\{|z|\ge241/240\}\cup\{\infty\}$ bounded by $C$, whose only marked
points are $v_a$ and $\infty$; in $\Delta\setminus\{\infty\}$, whose
fundamental group is $\Z$, the loop is homotopic to a simple loop once
around $\infty$, so \eqref{eq:twopunctures} gives $P(p\circ m)=T_C^{\pm1}$.
\end{proof}

The homotopy class of $\zeta$, not merely its homology class, must be
identified, since the isotopy class of $\gamma_a$ depends on it.

\begin{lemma}\label{lem:push}
The loop $\zeta$ is homotopic, as a loop based at $v_a$ in
$\Pone\setminus A_a$, to two clockwise circuits around the segment $[a,1]$
inside the disk bounded by $\eta$. Consequently
$h:=P(\zeta^{-1})$ satisfies $h(C)=T_\eta^{2}(C)=\gamma_a$, with the sign
of Convention~\ref{conv:sign}.
\end{lemma}

\begin{proof}
\emph{Step 1: enlarging the circle.} Put $R=\sqrt{1-a}=\sqrt{2/5}>a$. No
excluded point in \eqref{eq:forbidden} has modulus in $[a,R)$, so for
$a\le r<R$ the loops $\zeta_r(\theta)=Z(re^{2\pi i\theta},a)$ avoid $0,a,1,\infty$
and depend continuously on $r$. Their base points $Z(r,a)$ move along the
real path $\sigma_r=Z([a,r],a)\subset(1,v_a]$, since $Z(\cdot,a)$ is
decreasing on $(2/5,R)$ (the sign analysis in the proof of
Lemma~\ref{lem:m} is valid there) with $Z(R,a)=1$. Hence $\zeta$ is
based-homotopic in $\Pone\setminus A_a$ to $\sigma_r^{-1}\zeta_r\sigma_r$
for every $r\in[a,R)$.

\emph{Step 2: the circle $|x|=R$ maps onto $[a,1]$, and nothing inside it
does.} The substitution $u=x+(1-a)/x$ (a Joukowski map) simplifies $Z$:
writing $(x-1)^2-a=x(u-2)$ and $(x+a-1)(1-x)=x(2-a-u)$ gives
\begin{equation}\label{eq:Joukowski}
 Z(x,a)=G(u),\qquad G(u)=\frac{(u-2)^2}{4(2-a-u)}.
\end{equation}
For $x=Re^{i\vartheta}$ we have $u=2R\cos\vartheta$, so on $|x|=R$ the
variable $u$ runs over $[-2R,2R]$, where
$2-a-u\ge 2-a-2R>0$ and
\[
 G(u)-1=\frac{u^2-4(1-a)}{4(2-a-u)},\qquad
 G(u)-a=\frac{(u-2(1-a))^2}{4(2-a-u)} .
\]
Thus on $[-2R,2R]$ the function $G$ takes values in $[a,1]$, equals $1$
exactly at the endpoints, and equals $a$ exactly at $u=2(1-a)$; since
$G'(u)=(u-2)(2-2a-u)/(4(2-a-u)^2)$, $G$ decreases from $1$ to $a$ and then
increases back to $1$. Hence for $z\in(a,1)$ the equation $G(u)=z$ has two
distinct solutions $u\in(-2R,2R)$, and for each of them the two roots of
$x^2-ux+(1-a)=0$ are non-real conjugates with $|x|^2=1-a$, hence lie on
$|x|=R$. This produces four distinct points of $|x|=R$ mapping to $z$,
which exhaust the four preimages of $(a,z)$ under $p$. For the endpoints
the factorizations of $Z-1$ and $Z-a$ show that all preimages of $1$ and
of $a$ (namely $\pm R$ and $2/5\pm i\sqrt6/5$, with multiplicity) also lie
on $|x|=R$. So $Z(\{|x|=R\},a)=[a,1]$, and no $x$ with $|x|<R$ maps into
$[a,1]$.

\emph{Step 3: the homotopy class.} The segment $[a,1]$ lies in the open
disk bounded by $\eta$, so by Step 2 and continuity, for $r<R$ close to
$R$ the loop $\zeta_r$ lies in that disk and avoids $[a,1]$; the base
point path $\sigma_r\subset(1,v_a]$ also lies in the disk and avoids
$[a,1]$. So $\sigma_r^{-1}\zeta_r\sigma_r$ is a loop in the $\eta$-disk
minus the segment $[a,1]$. This space is an annulus, with fundamental
group $\Z$ generated by a counterclockwise circuit around $[a,1]$, and the
class of a loop is its winding number about any point of the segment, say
about $1$. Since the
homotopy of Step 1 avoids $1$, this winding number equals $\wind_1(\zeta)$,
which we compute by the argument principle applied to
\[
 Z(x,a)-1=-\frac{(x^2+a-1)^2}{4x(x-1)(x+a-1)}
\]
on the disk $|x|<a$: the numerator vanishes only at $\pm R$, outside the
disk, while the denominator has the two simple zeros $0$ and $1-a=2/5$
inside it. Hence $\wind_1(\zeta)=0-2=-2$, and $\zeta$ is based-homotopic to
two clockwise circuits around $[a,1]$ in the $\eta$-disk, as claimed.

\emph{Step 4: the point push.} Choose a round circle $\kappa$ enclosing
$a,1$ but not $v_a$ and lying inside both $C$ and $\eta$; for instance
$\kappa=\{|z-\frac45|=\frac{97}{480}\}$, for which
$|a-\frac45|=|1-\frac45|=\frac{96}{480}<\frac{97}{480}
<\frac{100}{480}=|v_a-\frac45|$, and the largest modulus on $\kappa$ is
$\frac{481}{480}<\frac{241}{240}$. The annulus between $\kappa$ and $\eta$
contains the single marked point $v_a$, and by Step 3, $\zeta^{-1}$ is
homotopic to two counterclockwise circuits of $v_a$ around this annulus.
By Convention~\ref{conv:sign} one such circuit has point push
$T_\eta T_\kappa^{-1}$, so, the twists about the disjoint curves $\kappa$
and $\eta$ commuting,
\[
 h=P(\zeta^{-1})=T_\eta^{2}T_\kappa^{-2}.
\]
Since $\kappa$ is disjoint from $C$, the twist $T_\kappa$ fixes the isotopy
class of $C$, and $h(C)=T_\eta^2(C)=\gamma_a$.
\end{proof}

\begin{proposition}\label{prop:g}
Let
\begin{equation}\label{eq:g}
 g=\ell^{-1}m\ell\in\pi_1(\W,b),
\end{equation}
the loop that traverses $\ell$, then $m$, then $\ell$ backwards. Then
$P(p\circ g)=T_{\gamma_a}^{\pm1}$: the deck transformation $g$ acts on
$\Teich(B_a)$ as the Dehn twist about $\gamma_a$ or its inverse. Both
coordinate loops of $g$ are null-homotopic in $Y$. In particular
$T_{\gamma_a}$ preserves $\D_{F_a}$.
\end{proposition}

\begin{proof}
Since $p\circ\ell$ moves only $v_a$, along $\zeta$, we have
$P(p\circ\ell)=P(\zeta)=h^{-1}$ with $h$ as in Lemma~\ref{lem:push}.
Multiplicativity of $P$ (Section~\ref{sec:conventions}) and
Lemma~\ref{lem:m} give
\[
 P(p\circ g)=P(p\circ\ell^{-1})\,P(p\circ m)\,P(p\circ\ell)
 =h\,T_C^{\pm1}\,h^{-1}=T_{h(C)}^{\pm1}=T_{\gamma_a}^{\pm1},
\]
using (P2) and Lemma~\ref{lem:push}. The $y$-coordinate of $g$ is constant.
The $x$-coordinate of $m$ is a small loop around $2/5$ together with a
real approach path in $(2/5,3/5]$; this loop lies in the disk
$|x-\frac12|<\frac12$, which avoids $0$ and $1$, so it is null-homotopic
in $Y$; hence the $x$-coordinate loop
of $g$ is homotopic to $(\pr_x\circ\ell)^{-1}(\pr_x\circ\ell)$, which is
trivial. Lemma~\ref{lem:criterion}(1) now shows that $g$, hence
$T_{\gamma_a}$, preserves $\D_{F_a}$.
\end{proof}

\section{No nonzero power of the twist preserves the base component}\label{sec:detector}

By Lemma~\ref{lem:criterion} and Proposition~\ref{prop:g}, it remains to
show that $g^n\notin E$ for $n\ne0$. Consider
\begin{equation}\label{eq:detector}
 \Phi\colon\W\longrightarrow Y\times\C^*,\qquad
 \Phi(x,y)=\bigl(t(x,y),L(x,y)\bigr)
 =\left(\frac{xy}{(1-x)(1-y)},\ (1-x)(1-y)\right).
\end{equation}
It is well defined: $L\ne0$ and $t\ne0$ since $x,y\in Y$, and
$t-1=(x+y-1)/((1-x)(1-y))\ne0$. In terms of the symmetric functions $e_1=x+y$, $e_2=xy$ one has
$L=1-e_1+e_2$ and $t=e_2/L$, so $\Phi$ is the symmetric quotient
$(x,y)\mapsto(x+y,xy)$ of Hironaka--Koch \cite[\S2.12]{HK}, restricted to
$\W$, up to this change of coordinates.

\emph{The constraint on diagonal loops.} Enlarge the diagonal by filling
in its six quadratic punctures:
$\widehat{\V}=\{(s,s):s\in\C\setminus\{0,1,1/2\}\}\supset\V$, on which
\eqref{eq:detector} still makes sense, so that the covering
\eqref{eq:square} below covers all of $Y$. In the coordinate $w=s/(1-s)$,
which maps $\C\setminus\{0,1,1/2\}$ bijectively onto
$\C\setminus\{0,1,-1\}$, the restriction of $\Phi$ is
\begin{equation}\label{eq:diag}
 t=w^2,\qquad L=(1+w)^{-2},\qquad w\in\C\setminus\{0,1,-1\},
\end{equation}
because $1+w=1/(1-s)$. The squaring map
\begin{equation}\label{eq:square}
 \C\setminus\{0,1,-1\}\longrightarrow Y,\qquad w\longmapsto w^2,
\end{equation}
is a two-sheeted covering; our base points are $w_0=a/(1-a)=3/2$ and
$t_0=w_0^2=9/4$.

\begin{lemma}\label{lem:constraint}
Let $\lambda$ be a loop in $\V$ based at $b$. Then $w\circ\lambda$ is the
lift of the loop $t\circ\lambda$ through \eqref{eq:square} starting at
$w_0$; it is closed, and
\begin{equation}\label{eq:constraint}
 \wind_0(L\circ\lambda)=-2\,\wind_{-1}(w\circ\lambda).
\end{equation}
The based homotopy class of $w\circ\lambda$ in $\C\setminus\{0,1,-1\}$,
hence the right-hand side of \eqref{eq:constraint}, depends only on the
based homotopy class of $t\circ\lambda$ in $Y$.
\end{lemma}

\begin{proof}
On $\V\subset\widehat\V$ formula \eqref{eq:diag} applies: $w\circ\lambda$
is a closed path in $\C\setminus\{0,1,-1\}$ from $w_0$ with
$(w\circ\lambda)^2=t\circ\lambda$, hence the lift of $t\circ\lambda$ from
$w_0$, and $L\circ\lambda=(1+w\circ\lambda)^{-2}$ gives
\eqref{eq:constraint}. The last assertion is homotopy lifting for the
covering \eqref{eq:square}.
\end{proof}

\emph{The loops $g^n$ violate the constraint.}
On the fibre $y=a$ the map \eqref{eq:detector} becomes
\begin{equation}\label{eq:fiber}
 t=\frac{3x}{2(1-x)},\qquad L=\frac25(1-x).
\end{equation}
Let $\alpha=t\circ\ell$ and $\beta=t\circ m$, loops in $Y$ based at $t_0$.

\begin{lemma}\label{lem:gn}
For every integer $n$, the loop $g^n$ satisfies
\begin{equation}\label{eq:mismatch}
 \wind_0(L\circ g^n)=0,\qquad
 \wind_{-1}\bigl(\widetilde{t\circ g^n}\bigr)=n,
\end{equation}
where $\widetilde{t\circ g^n}$ is the lift of $t\circ g^n$ through
\eqref{eq:square} starting at $w_0$, which is a closed loop.
\end{lemma}

\begin{proof}
\emph{The loop $\ell$.} Along $\ell$ we have $x=ae^{2\pi i\theta}$ with
$|x|=a<1$, so $1-x$ does not wind about $0$ and \eqref{eq:fiber} gives
$\wind_0(\alpha)=1$ and $\wind_0(L\circ\ell)=0$. Since $\wind_0(\alpha)$
is odd, continuation of the square root along
$\alpha$, or along $\alpha^{-1}$, interchanges the two sheets of
\eqref{eq:square}: the lift $\tilde\alpha$ of $\alpha$ starting at $w_0$
ends at $-w_0$.

\emph{The loop $m$.} The M\"obius map $x\mapsto t$ in \eqref{eq:fiber}
sends $2/5$ to $1$ and $a$ to $9/4$, is increasing on the real interval
$(2/5,a]$, which it maps onto $(1,9/4]$, and, being conformal, sends a
small counterclockwise circle about $2/5$ to a small counterclockwise loop
about $1$. Hence $\beta$ is a positive meridian about $1$ in $Y$,
approached from $t_0$ along the real axis to the right of $1$. Its lift
starting at $w_0=3/2$ runs along $(1,3/2]$ and around a small loop about
$+1$, so it is a positive meridian about $+1$; the lift starting at $-w_0$
is its negative, a positive meridian about $-1$ (the map $w\mapsto-w$
preserves orientation). Also $\wind_0(L\circ m)=0$, since $1-x$ stays near
the positive interval $[2/5,3/5]$ along $m$.

\emph{The loop $g$.} By additivity,
$\wind_0(L\circ g)=-\wind_0(L\circ\ell)+\wind_0(L\circ m)+\wind_0(L\circ\ell)=0$,
and so $\wind_0(L\circ g^n)=0$. The $t$-projection of $g$ is
$\alpha^{-1}\beta\alpha$: first $\alpha$, then $\beta$, then $\alpha$
backwards. Its lift from $w_0$ consists of $\tilde\alpha$, ending on the
other sheet at $-w_0$, then the lift of $\beta$ from $-w_0$, a positive
meridian about $-1$, then $\tilde\alpha$ reversed, returning to $w_0$. The
contributions of $\tilde\alpha$ and its reverse to the change of argument
of $w+1$ cancel, so this closed lift has winding number $1$ about $-1$,
and the lift of $t\circ g^n=(\alpha^{-1}\beta\alpha)^n$ from $w_0$, the
$n$-fold concatenation, has winding number $n$.
\end{proof}

\begin{proposition}\label{prop:notinE}
For every $n\ne0$, $g^n\notin E$. Consequently Theorem~\ref{thm:main}
holds for the model map $F_a$ and the curve $\gamma_a=T_\eta^2(C)$.
\end{proposition}

\begin{proof}
Suppose $g^n\in E$, so that $g^n$ is homotopic in $\W$ to a loop $\lambda$
in $\V$ based at $b$. Applying $\Phi$, the loops $t\circ\lambda$ and
$t\circ g^n$ are based-homotopic in $Y$, and $L\circ\lambda$ and
$L\circ g^n$ are homotopic in $\C^*$. By Lemma~\ref{lem:constraint}, the
lift $w\circ\lambda$ is based-homotopic to the lift of $t\circ g^n$ used in
Lemma~\ref{lem:gn}, so $\wind_{-1}(w\circ\lambda)=n$, while
$\wind_0(L\circ\lambda)=\wind_0(L\circ g^n)=0$. The constraint
\eqref{eq:constraint} then reads $0=-2n$, so $n=0$.

For the second assertion, $g$ acts as $T_{\gamma_a}^{\pm1}$ and has
homotopic coordinate loops (Proposition~\ref{prop:g}), so by
Lemma~\ref{lem:criterion} the twist $T_{\gamma_a}$ preserves $\D_{F_a}$
and the components $T_{\gamma_a}^{n}(D_0)=g^{\pm n}(D_0)$, $n\in\Z$, are
pairwise distinct.
\end{proof}

\begin{remark}\label{rem:othersign}
The same computation applies to the other conjugate $g'=\ell m\ell^{-1}$
(its $t$-projection $\alpha\beta\alpha^{-1}$ also changes sheet before the
meridian, as $\alpha^{-1}$ has odd winding number about $0$, and again
$\wind_0(L\circ g')=0$), whose mapping class is
$h^{-1}T_C^{\pm1}h=T_{h^{-1}(C)}^{\pm1}$ with
$h^{-1}(C)=T_\eta^{-2}T_\kappa^{2}(C)=T_\eta^{-2}(C)$. Hence
Theorem~\ref{thm:main} holds for both curves $T_\eta^{\pm2}(C)$.
\end{remark}

\section{Transport to an arbitrary map of the family}\label{sec:transport}

Let $f$, $\mu$, $F_s$, $s_t$, $H_t$ be as in Construction~\ref{con:transport}.
We show that changing markings by $H=H_1$ carries $\D_{F_a}$ onto
$\D_{F_s}$, $D_0$ to the base component of $\D_{F_s}$, and conjugates
$T_{\gamma_a}$ to $T_{H(\gamma_a)}$.

\begin{lemma}\label{lem:lift}
There is a continuous family of homeomorphisms $K_t$ of $\Pone$,
$0\le t\le1$, with $K_0=\id$ and
\begin{equation}\label{eq:equiv}
 F_{s_t}\circ K_t=H_t\circ F_a\qquad(0\le t\le1).
\end{equation}
It satisfies $K_t|_{A_a}=H_t|_{A_a}$ for all $t$; consequently
$H_t^{-1}K_t$ is an isotopy from the identity through homeomorphisms
fixing $A_a$ pointwise, and $K_1$ is isotopic to $H_1$ relative to $A_a$.
\end{lemma}

\begin{proof}
Every $F_s$ is M\"obius-equivalent to $Q(z)=z^2$, with M\"obius maps
depending continuously on $s$: completing the square in
$F_s(z)=1+b_s/z+c_s/z^2$ and recalling $v(s)=1-b_s^2/(4c_s)$ gives
$F_s=v(s)+c_s\,J_s^2$, that is,
\[
 \nu_s\circ F_s=Q\circ J_s,\qquad
 J_s(z)=\frac1z+\frac{b_s}{2c_s},\qquad \nu_s(w)=\frac{w-v(s)}{c_s},
\]
with $c_s\ne0$ for $s\in U$. The isotopy
$\widehat H_t=\nu_{s_t}\circ H_t\circ \nu_a^{-1}$ starts at the identity
and fixes $0$ and $\infty$, because $H_t(v_a)=v(s_t)$ and
$H_t(\infty)=\infty$. By the homotopy lifting property of the covering
$Q\colon\C^*\to\C^*$, the homotopy $(z,t)\mapsto\widehat H_t(Q(z))$ lifts
uniquely to a continuous $\rho\colon\C^*\times[0,1]\to\C^*$ with
$\rho_0=\id$ and $Q\circ\rho_t=\widehat H_t\circ Q$; lifting
$\widehat H_t^{-1}\circ Q$ in the same way gives $\rho'_t$ with
$\rho'_t\rho_t=\id=\rho_t\rho'_t$ (both products lift the constant
homotopy $Q$ and start at $\id$), so each $\rho_t$ is a homeomorphism of
$\C^*$, and $\rho$ extends continuously to $\Pone\times[0,1]$ by
$\rho_t(0)=0$, $\rho_t(\infty)=\infty$, since
$|\rho_t(z)|^2=|\widehat H_t(z^2)|$ and $\widehat H$ is continuous on the
compact space $\Pone\times[0,1]$. Then
$K_t=J_{s_t}^{-1}\circ\rho_t\circ J_a$ is a continuous family of
homeomorphisms with $K_0=\id$ and
$F_{s_t}K_t=\nu_{s_t}^{-1}Q\rho_tJ_a=\nu_{s_t}^{-1}\widehat H_tQJ_a
=H_t\nu_a^{-1}QJ_a=H_tF_a$, which is \eqref{eq:equiv}.

For the marked points, \eqref{eq:equiv} gives
$F_{s_t}(K_t(w))=H_t(F_a(w))$ for $w\in A_a$. Thus $K_t(0)$ is the critical
point of $F_{s_t}$ over $\infty$, namely $0$; $K_t(\infty)$ lies in
$F_{s_t}^{-1}(1)=\{\infty,-c_{s_t}/b_{s_t}\}$, $K_t(1)$ in
$F_{s_t}^{-1}(s_t)$, and $K_t(a)$ in $F_{s_t}^{-1}(0)\ni s_t$. For $s_t\in U$ each of the latter three fibres
(over the non-critical values $1$, $s_t$, $0$) consists of
two distinct points depending continuously on $t$: $b_{s_t}\ne0$ since
$s_t^2+s_t-1\ne0$, and $s_t$, $0$ are not critical values of $F_{s_t}$
since $v(s_t)\ne s_t,0$. By continuity from $t=0$,
$K_t(0)=0$, $K_t(\infty)=\infty$, $K_t(1)=1$ and $K_t(a)=s_t$, which are
the values of $H_t$ on $A_a$. The remaining assertions follow.
\end{proof}

\begin{proof}[Proof of Theorem~\ref{thm:main}]
Write $H=H_1$, $K=K_1$. Since $H$ maps $B_a$ onto $B_s$ (notation of
Section~\ref{sec:curve}), changing markings by $H$ is a biholomorphic
isomorphism
\[
 \Psi\colon\Teich(B_a)\longrightarrow\Teich(B_s),\qquad
 [\phi]\longmapsto[\phi\circ H^{-1}],
\]
and likewise $[\psi]\mapsto[\psi\circ H^{-1}]$ identifies $\Teich(A_a)$
with $\Teich(A_s)$, $A_s=\{0,\infty,1,s\}=H(A_a)=K(A_a)$.

\emph{$\Psi$ carries $\D_{F_a}$ onto $\D_{F_s}$.} Let $\phi$ be a marking
with pullback diagram $\widetilde F\circ\psi=\phi\circ F_a$ for $F_a$.
Using $F_a\circ K^{-1}=H^{-1}\circ F_s$ from \eqref{eq:equiv},
\[
 \widetilde F\circ(\psi K^{-1})=\phi\circ F_a\circ K^{-1}
 =(\phi H^{-1})\circ F_s ,
\]
so $\psi K^{-1}$ is a source marking for the target marking $\phi H^{-1}$
in a pullback diagram for $F_s$: $F_s^*\Psi[\phi]=[\psi K^{-1}]$, while
$i_*\Psi[\phi]=[\phi H^{-1}]$. Now $[\psi K^{-1}]=[\phi H^{-1}]$ in
$\Teich(A_s)$ if and only if $[\psi K^{-1}H]=[\phi]$ in $\Teich(A_a)$, and
$K^{-1}H$ is isotopic to the identity relative to $A_a$ by
Lemma~\ref{lem:lift}, so this holds if and only if $[\psi]=[\phi]$ in
$\Teich(A_a)$, that is, $[\phi]\in\D_{F_a}$. Hence
$\Psi(\D_{F_a})=\D_{F_s}$.

\emph{$\Psi$ conjugates twists.} For a mapping class $k\in\PMod(B_a)$,
$\Psi([\phi]\cdot k)=[\phi kH^{-1}]=[\phi H^{-1}]\cdot(HkH^{-1})$, so $\Psi$
intertwines the action of $T_{\gamma_a}$ with that of
$HT_{\gamma_a}H^{-1}=T_{H(\gamma_a)}$, by the naturality of twists (P2),
which holds verbatim for the homeomorphism $H\colon(\Pone,B_a)\to(\Pone,B_s)$.

\emph{$\Psi$ carries $D_0$ to the base component.} The path of markings
$[H_t]$ in $\Teich(B_a)$ lies in $\D_{F_a}$: by \eqref{eq:equiv}, $K_t$ is
a source marking for the target marking $H_t$, and $[K_t]=[H_t]$ in
$\Teich(A_a)$ because $H_t^{-1}K_t$ is isotopic to the identity relative to
$A_a$ (Lemma~\ref{lem:lift}). The path starts at $\tau_0$, so its endpoint
$[H]$ lies in $D_0$, and $\Psi[H]=[\id]$ is the base point of
$\Teich(B_s)$. Hence $\Psi(D_0)$ is the base component of $\D_{F_s}$, and
the components $T_{H(\gamma_a)}^n(\Psi(D_0))=\Psi(T_{\gamma_a}^n(D_0))$,
$n\in\Z$, are pairwise distinct by Proposition~\ref{prop:notinE}. Finally,
conjugating by $\mu^{-1}$ identifies $(\Pone,B_s,F_s)$ with $(\Pone,B,f)$
and $H(\gamma_a)$ with $\gamma$, which proves the theorem for $f$.
\end{proof}

\begin{thebibliography}{KPS}
\bibitem[BCT]{BCT}
X. Buff, G. Cui and Tan Lei,
\emph{Teichm\"uller spaces and holomorphic dynamics},
in: Handbook of Teichm\"uller Theory, Vol.\ IV, IRMA Lect.\ Math.\ Theor.\
Phys.\ 19, Eur.\ Math.\ Soc., 2014, 717--756.
\bibitem[Ep]{Ep}
A. L. Epstein,
\emph{Transversality in holomorphic dynamics},
manuscript.
\bibitem[FM]{FM}
B. Farb and D. Margalit,
\emph{A Primer on Mapping Class Groups},
Princeton Math.\ Ser.\ 49, Princeton Univ.\ Press, 2012.
\bibitem[FKS]{FKS}
T. Firsova, J. Kahn and N. Selinger,
\emph{On deformation spaces of quadratic rational functions},
Int.\ Math.\ Res.\ Not.\ IMRN 2023, no.\ 8, 6703--6738; arXiv:1606.00723v4
(2020).
\bibitem[HK]{HK}
E. Hironaka and S. Koch,
\emph{A disconnected deformation space of rational maps},
J.\ Mod.\ Dyn.\ 11 (2017), 409--423.
\bibitem[HKa]{HKa}
E. Hironaka and S. Koch,
\emph{A disconnected deformation space of rational maps},
arXiv:1602.07378v1 (2016).
\bibitem[Ko]{Ko}
S. Koch,
\emph{Teichm\"uller theory and critically finite endomorphisms},
Adv.\ Math.\ 248 (2013), 573--617.
\bibitem[KPS]{KPS}
S. Koch, K. M. Pilgrim and N. Selinger,
\emph{Pullback invariants of Thurston maps},
Trans.\ Amer.\ Math.\ Soc.\ 368 (2016), 4621--4655.
\end{thebibliography}

\end{document}
